\addtocontents{toc}{\protect\setcounter{tocdepth}{-1}} % Schaltet das TOC kurzzeitig "stumm"
\chapter{Einleitung}\label{ch:einleitung}             % Erzeugt "1 Einleitung" im Text
\addtocontents{toc}{\protect\setcounter{tocdepth}{2}}

Die globale Zunahme von Waldbrandereignissen hat sich in den letzten Jahren von einem saisonalen Phänomen zu einer chronischen, weltweiten Bedrohung entwickelt. Die ökologischen, sozialen und ökonomischen Auswirkungen von Waldbränden sind verheerend: Neben der Zerstörung sensibler Ökosysteme und dem Verlust von Biodiversität führen Waldbrände zur massiven Freisetzung von Treibhausgasen, was die anthropogenen Erderwärmung weiter verstärkt.
Die Tatsache, dass im Jahr 2024 erstmals die \qty{1,5}{\degree}C-Grenze überschritten wurde, verdeutlicht die Dringlichkeit, denn die zunehmende Austrocknung von Böden und Vegetation begünstigt die Entstehung von Feuern erheblich~\cite{Copernicus2025,Umweltbundesamt2024}.
Längst beschränken sich die wirtschaftlichen Folgen nicht mehr nur auf Regionen wie Kalifornien, wo die staatlichen Budgets bereits im Jahr 2020 um 3,1 Milliarden USD übertroffen wurden~\cite{LegislativeAnalystsOffice2020}. Parallel dazu sind auch in Europa die jährlichen Ausgaben für die Brandbekämpfung auf inzwischen rund zwei Milliarden Euro angestiegen~\cite{JointResearchCentre2023}.

Die entscheidende Variable zur Eindämmung dieser Schäden ist die Zeit. Die Phasen zwischen dem Brandausbruch und der offenen Flamme, die sogenannte Schwelbrandphase, bietet das größte Potenzial für eine effektive Brandbekämpfung mit minimalem Ressourceneinsatz~\cite{TobisJakobs}. Herkömmliche Detektionsmethoden stoßen hier an ihre Grenzen. Während Meldungen durch die Bevölkerung von der lokalen Siedlungsdichte abhängig sind, sind bemannte Feuerwachtürme aufgrund von mangelnder Präzision und arbeitsrechtlicher Bestimmungen kaum noch zukunftsfähig~\cite{DirkSchneider2017}.

Die technische Antwort auf dieses Defizit ist die Implementierung moderner Früherkennungssysteme, die jedoch spezifische Schwachstellen aufweisen. Weltraumgestützte Systeme leiden unter niedrigen Wiederholraten, was eine Echtzeit-Warnung erschwert. Optisch-terrestrische Kamerasysteme sind schneller, reagieren jedoch auf topografische Hindernisse oder Wetterphänomene, was häufig zu ressourcenintensiven Fehlalarmen führt~\cite{Chagger2018}. Internet of Things (\acs{IoT})-basierte Sensoren wiederum versprechen eine präzise Detektion, erfordern jedoch für eine flächendeckende Abdeckung eine kostspielige und wartungsintensive Infrastruktur~\cite{Klick2021}.

Hieraus ergibt sich eine kritische Diskrepanz zwischen den theoretischen Potenzialen einzelner Technologien und den praktischen Anforderungen an ein zuverlässiges System, das eine geringe Latenzzeit mit einer niedrigen Fehlalarmrate verbindet. Ziel der vorliegenden Arbeit ist es daher, diese technologische Lücke durch eine vergleichende Evaluierung zu analysieren. Im Fokus steht die Frage, welche Kombination aus sensorischen und optischen Frühwarnsystemen unter Berücksichtigung von Latenzzeit, Fehlalarmrate und Kosteneffizienz den effektivsten Schutz bietet.

Es wird die Hypothese aufgestellt, dass multimodale Systeme, die lokale chemische Sensoren mit optischen Systemen kombinieren und dabei weltraumgestützte Frühwarnsysteme nutzen, eine signifikant höhere Zuverlässigkeit und Geschwindigkeit bieten als isolierte Einzellösungen.